%% explicacion_plan_de_vuelo.tex
%% Derivación didáctica del optimizador del plan de vuelo del receptor
%% en el sistema SAR biestático helicoidal con dos medios homogéneos.
%%
%% Compilar: pdflatex explicacion_plan_de_vuelo.tex (dos pasadas)
%% Figuras:  hypotheses/figures/pv_fig0X_*.pdf (generar con utils/generate_figures_plan_vuelo.py)

\documentclass[12pt,a4paper,oneside]{book}

% ─── paquetes ────────────────────────────────────────────────────────────────
\usepackage[utf8]{inputenc}
\usepackage[T1]{fontenc}
\usepackage[spanish]{babel}
\usepackage{lmodern}
\usepackage{geometry}
\geometry{left=3cm, right=2.5cm, top=2.5cm, bottom=2.5cm}

\usepackage{amsmath,amssymb,amsfonts,amsthm}
\usepackage{mathtools}
\usepackage{bm}
\usepackage{physics}
\usepackage{siunitx}
\sisetup{output-decimal-marker={,}}

\usepackage{graphicx}
\graphicspath{{figures/}}

\usepackage{float}
\usepackage{caption}
\usepackage{subcaption}
\usepackage[table,xcdraw]{xcolor}
\usepackage{booktabs}
\usepackage{array}
\usepackage{multirow}
\usepackage{algorithm}
\usepackage{algpseudocode}

\usepackage{tikz}
\usetikzlibrary{arrows.meta,decorations.pathreplacing,calc,babel,patterns}

\usepackage{mdframed}
\usepackage{tcolorbox}
\tcbuselibrary{theorems,skins,breakable}

\usepackage{hyperref}
\hypersetup{
  colorlinks=true,
  linkcolor=blue!60!black,
  citecolor=green!50!black,
  urlcolor=cyan!60!black,
  pdftitle={Optimización del Plan de Vuelo en SAR Biestático},
  pdfauthor={Investigación SAR-Tomo},
}

\usepackage{fancyhdr}
\setlength{\headheight}{15pt}
\pagestyle{fancy}
\fancyhf{}
\fancyhead[L]{\small\nouppercase{\leftmark}}
\fancyhead[R]{\small\thepage}
\renewcommand{\headrulewidth}{0.4pt}

% ─── entornos ─────────────────────────────────────────────────────────────────
\newcommand{\result}[2]{%
  \begin{tcolorbox}[enhanced, breakable,
    colback=yellow!8, colframe=orange!70!black,
    fonttitle=\bfseries,
    title={\protect#1},
    attach boxed title to top left={yshift=-2mm, xshift=5mm},
    boxed title style={colback=orange!70!black}]
    #2
  \end{tcolorbox}}

\newcommand{\defin}[2]{%
  \begin{tcolorbox}[enhanced, breakable,
    colback=blue!5, colframe=blue!50!black,
    fonttitle=\bfseries,
    title={\protect#1}]
    #2
  \end{tcolorbox}}

\newcommand{\nota}[2]{%
  \begin{tcolorbox}[enhanced, breakable,
    colback=green!5, colframe=green!50!black,
    fonttitle=\bfseries\small,
    title={\protect#1}]
    #2
  \end{tcolorbox}}

\newcommand{\atencion}[2]{%
  \begin{tcolorbox}[enhanced, breakable,
    colback=red!5, colframe=red!60!black,
    fonttitle=\bfseries\small,
    title={\protect#1}]
    #2
  \end{tcolorbox}}

% ─── macros matemáticas ───────────────────────────────────────────────────────
\newcommand{\RT}{R_T}
\newcommand{\RR}{R_R}
\newcommand{\thetaB}{\theta_B}
\newcommand{\nuno}{n_1}
\newcommand{\ndos}{n_2}
\newcommand{\Prcv}{P_r}
\newcommand{\Pt}{P_t}
\newcommand{\Jcost}{J}
\newcommand{\alphaR}{\alpha}
\newcommand{\dxy}{\delta_{xy}}
\newcommand{\dz}{\delta_z}

% ─────────────────────────────────────────────────────────────────────────────
\begin{document}

\frontmatter

% ─── portada ──────────────────────────────────────────────────────────────────
\begin{titlepage}
\centering
\vspace*{2cm}
{\LARGE\bfseries Optimización del Plan de Vuelo del Receptor}\\[0.5cm]
{\LARGE\bfseries en SAR Biestático Helicoidal}\\[0.3cm]
{\LARGE\bfseries Ángulo de Brewster, Ecuación de Radar y}\\[0.3cm]
{\LARGE\bfseries Resolución 3D Multiobjetivo}\\[2cm]
{\large\itshape Derivación paso a paso: desde la física de Fresnel\\
hasta la función de costo log-normalizada y\\
la estrategia de búsqueda multi-inicio}\\[3cm]
{\large Proyecto de Investigación SAR-Tomo}\\[0.5cm]
{\large\today}\\[4cm]
\includegraphics[width=0.65\textwidth]{pv_fig10_sistema_3d.pdf}\\[1cm]
{\small Configuración biestática óptima: TX en hélice cónica (azul), RX
optimizado para ángulo de Brewster (rojo), targets subsuperficiales (×),
y caminos refractados en la interfaz aire-suelo ($z=0$).}
\end{titlepage}

\tableofcontents

\mainmatter

%%═══════════════════════════════════════════════════════════════════════════════
\chapter{Introducción: El Problema del Plan de Vuelo}
%%═══════════════════════════════════════════════════════════════════════════════

\section{Motivación}

En un sistema SAR biestático para tomografía de subsuelo, el transmisor (TX) sigue una trayectoria predeterminada —la hélice cónica— optimizada para maximizar la diversidad angular y la resolución 3D del volumen de interés. Sin embargo, la posición del receptor (RX) en cada instante es un \textbf{grado de libertad libre} que el diseñador del sistema puede elegir.

La pregunta central de este documento es:

\begin{center}
\fbox{\parbox{0.82\textwidth}{\centering\itshape
¿Dónde debe volar el receptor en cada instante para que el sistema biestático\\
sea simultáneamente óptimo en potencia recibida \textbf{y} en resolución 3D?
}}
\end{center}

Esta pregunta, aparentemente simple, tiene varias capas de complejidad:

\begin{itemize}
  \item La potencia recibida depende de la geometría biestática a través de la \textbf{ecuación de radar} y los \textbf{coeficientes de Fresnel} en la interfaz aire-suelo.
  \item La resolución 3D depende del \textbf{ángulo de separación biestática} $\Delta\phi$ y del ángulo de observación $\psi_0$, como se derivó en \texttt{explicacion\_resolucion.tex}.
  \item Ambos criterios tienen \textbf{escalas y dimensiones completamente distintas} (Watts vs. metros), lo que hace imposible sumarlos directamente.
  \item El espacio de búsqueda es \textbf{no convexo}: múltiples mínimos locales en la función de costo.
\end{itemize}

\section{Enfoque adoptado}

El optimizador del plan de vuelo (\texttt{run\_plano\_de\_voo.m}) resuelve este problema con tres ideas clave:

\begin{enumerate}
  \item \textbf{Explotación del ángulo de Brewster}: la polarización TM incidente bajo un ángulo específico produce transmitancia unitaria en la interfaz, maximizando la potencia que llega al target subsuperficial.
  \item \textbf{Normalización logarítmica}: la función de costo combina potencia y resolución usando $\log_{10}$, haciendo ambos términos adimensionales y de magnitud comparable.
  \item \textbf{Multi-inicio guiado geométricamente}: la búsqueda con \texttt{fmincon} parte de 5 candidatos heurísticos que cubren las regiones de interés biestático.
\end{enumerate}

\section{Relación con otros documentos}

Este documento se apoya directamente en:
\begin{itemize}
  \item \texttt{explicacion\_resolucion.tex}: las fórmulas de resolución 3D biestática ($\dxy$, $\dz$) que se usan como segundo objetivo de optimización.
  \item Los archivos MATLAB en \texttt{simulations/flightpath/}: implementación numérica de todos los conceptos derivados aquí.
\end{itemize}

%%═══════════════════════════════════════════════════════════════════════════════
\chapter{Física de la Interfaz: Polarización TM y Ángulo de Brewster}
%%═══════════════════════════════════════════════════════════════════════════════

\section{La interfaz aire-suelo}

El sistema opera sobre dos medios homogéneos separados por una interfaz plana horizontal en $z=0$:
\begin{itemize}
  \item \textbf{Medio 1 (aire)}: $z > 0$, índice de refracción $\nuno = 1$, velocidad $c$.
  \item \textbf{Medio 2 (suelo)}: $z < 0$, índice de refracción $\ndos = \sqrt{\varepsilon_r}$, velocidad $v_2 = c/\ndos$.
\end{itemize}

Para el suelo seco típico, $\varepsilon_r \approx 4$ $\Rightarrow$ $\ndos = 2$.

La señal radar de banda P ($f_0 = 400\,\text{MHz}$) incide oblicuamente en la interfaz con un ángulo $\theta_1$ (medido desde la normal). Una fracción de la potencia se transmite al medio 2 con ángulo $\theta_2$ dado por la ley de Snell:
\[
\nuno\sin\theta_1 = \ndos\sin\theta_2
\quad\Longrightarrow\quad
\sin\theta_2 = \frac{\sin\theta_1}{\ndos}
\]

\section{Las ecuaciones de Fresnel para polarización TM}

La señal radar utiliza polarización \textbf{TM} (campo eléctrico en el plano de incidencia, también llamada polarización p). Las ecuaciones de Fresnel para los coeficientes de campo eléctrico son:

\result{Coeficientes de Fresnel — Polarización TM}{
\begin{align}
r_{TM} &= \frac{\ndos\cos\theta_1 - \nuno\cos\theta_2}{\ndos\cos\theta_1 + \nuno\cos\theta_2}
\label{eq:r_TM}\\[6pt]
t_{TM} &= \frac{2\nuno\cos\theta_1}{\ndos\cos\theta_1 + \nuno\cos\theta_2}
\label{eq:t_TM}
\end{align}
}

Para la ecuación de radar, lo relevante es la \textbf{transmitancia y reflectancia de potencia}, que incluyen el factor de cambio de impedancia entre medios:

\result{Transmitancia y Reflectancia de Potencia}{
\begin{align}
T &= \frac{\ndos\cos\theta_2}{\nuno\cos\theta_1}\,|t_{TM}|^2
\label{eq:T_potencia}\\[6pt]
R &= |r_{TM}|^2
\label{eq:R_potencia}
\end{align}
con la conservación de energía $T + R = 1$ $\checkmark$
}

\section{El ángulo de Brewster: transmitancia unitaria}

Para polarización TM existe un ángulo especial en el que el coeficiente de reflexión se anula exactamente: $r_{TM} = 0$. Esto ocurre cuando numerador de~\eqref{eq:r_TM} es cero:
\[
\ndos\cos\theta_1 = \nuno\cos\theta_2
\]
Combinando con la ley de Snell $\sin\theta_2 = \sin\theta_1/\ndos$:

\result{Ángulo de Brewster}{
\[
\boxed{\thetaB = \arctan\!\left(\frac{\ndos}{\nuno}\right)}
\]
En $\theta_1 = \thetaB$: $r_{TM} = 0$, $T = 1$ (toda la potencia se transmite al subsuelo).
Para $\ndos = 2$: $\thetaB = \arctan(2) \approx 63.43^\circ$.
}

La figura~\ref{fig:fresnel_tm} muestra las curvas $T(\theta_1)$ y $R(\theta_1)$ completas.

\begin{figure}[H]
\centering
\includegraphics[width=0.85\textwidth]{pv_fig01_fresnel_tm.pdf}
\caption{Coeficientes de Fresnel para polarización TM con $\ndos=2$ ($\varepsilon_r=4$).
(a) Transmitancia $T$ y reflectancia $R$ de potencia en función del ángulo de incidencia.
(b) Coeficientes de campo $|r_{TM}|$ y $|t_{TM}|$.
El ángulo de Brewster $\thetaB \approx 63.43^\circ$ es el único punto donde $R=0$, $T=1$.
Para el sistema P-band, este ángulo es el objetivo de la optimización del plan de vuelo.}
\label{fig:fresnel_tm}
\end{figure}

\section{Consecuencia para la geometría biestática}

Si el RX se coloca de modo que la señal llega al target con ángulo de incidencia $\theta_1 \approx \thetaB$ \emph{tanto en el tramo TX→Target como en el tramo Target→RX}, la potencia transmitida es máxima. Esta condición impone restricciones en la posición horizontal del RX relativa al TX y al centroide de los targets.

\nota{Cuantificación de la ganancia de Brewster}{
Para $\ndos=2$ y ángulo de incidencia de $63.43^\circ$: $T=1.0$ (Brewster). Para ángulo de $30^\circ$: $T\approx 0.88$. Para incidencia normal ($0^\circ$): $T\approx 0.89$. El ángulo de Brewster proporciona $\sim 12-14\%$ más de potencia transmitida respecto al promedio, lo cual es modesto pero real.
}

La figura~\ref{fig:brewster_geometry} muestra la geometría óptima que sitúa la señal biestática en el ángulo de Brewster.

\begin{figure}[H]
\centering
\includegraphics[width=0.80\textwidth]{pv_fig08_brewster_geometry.pdf}
\caption{Geometría óptima de Brewster en el plano biestático. El transmisor TX ilumina
el target con ángulo de incidencia $\thetaB$ sobre la interfaz. El receptor RX se ubica
en la posición simétrica (opuesta al TX respecto al target) para que el tramo Target→RX
también incida a $\thetaB$. La \emph{distancia de Brewster} $d_B = (z_{RX}+|z_P|)/\tan\thetaB$
define el radio horizontal óptimo del RX desde el target.}
\label{fig:brewster_geometry}
\end{figure}

%%═══════════════════════════════════════════════════════════════════════════════
\chapter{El Principio de Fermat y la Solución Analítica de Snell}
%%═══════════════════════════════════════════════════════════════════════════════

\section{El principio de Fermat}

\defin{Principio de Fermat (formulación de Hamilton)}{
La trayectoria seguida por un rayo electromagnético entre dos puntos es aquella que hace
\emph{estacionario} (en general, mínimo) el tiempo de tránsito óptico:
\[
\mathcal{T} = \int_A^B n(\mathbf{r})\,ds
\]
donde $n(\mathbf{r})$ es el índice de refracción local y $ds$ el diferencial de longitud de trayectoria.
}

Para dos medios homogéneos con interfaz plana, el tiempo óptico entre fuente $S=(s_x,s_y,s_z)$ con $s_z>0$ y target $T=(t_x,t_y,t_z)$ con $t_z<0$ a través de un punto de refracción $P=(p_x,p_y,0)$ en la interfaz es:
\[
\mathcal{T}(P) = \nuno\,\|S-P\| + \ndos\,\|P-T\|
= \nuno\,d_1(P) + \ndos\,d_2(P)
\]

Minimizar $\mathcal{T}$ sobre $P$ da exactamente la \textbf{ley de Snell} en el punto de refracción.

\section{Reducción 3D $\rightarrow$ 1D para interfaz plana horizontal}

Un resultado clave, que permite evitar el optimizador 2D completo, es la siguiente propiedad geométrica:

\result{Reducción Dimensional de Fermat}{
Para una interfaz plana horizontal ($z=0$), el punto de refracción óptimo $P^*$ siempre pertenece
a la \textbf{línea recta horizontal} que une las proyecciones de la fuente y del target sobre
la interfaz:
\[
P^* = S_{xy} + u^*\,\hat{d}_{xy}, \quad u^* \in [0, D]
\]
donde $S_{xy}=(s_x,s_y)$, $T_{xy}=(t_x,t_y)$, $\hat{d}_{xy} = (T_{xy}-S_{xy})/D$ y $D=\|T_{xy}-S_{xy}\|$.

En consecuencia, el problema 2D (buscar $p_x,p_y$) se reduce a un problema 1D (encontrar el escalar $u^*$).
}

\textbf{Demostración:} La condición de primer orden $\nabla_P\,\mathcal{T}=0$ en el plano $z=0$ requiere:
\[
\nuno\frac{P-S_{xy}}{d_1} + \ndos\frac{P-T_{xy}}{d_2} = \mathbf{0}
\]
Esto significa que $(P-S_{xy})$ y $(T_{xy}-P)$ son paralelos (colineales): $P$ debe estar en la línea $S_{xy}$-$T_{xy}$. $\square$

\section{La ecuación de Snell 1D y su solución Newton-Raphson}

Parametrizando $P = S_{xy} + u\,\hat{d}_{xy}$ con $u\in[0,D]$:
\[
d_1(u) = \sqrt{u^2 + s_z^2}, \qquad d_2(u) = \sqrt{(D-u)^2 + |t_z|^2}
\]

La condición de Snell $\nuno\sin\theta_1 = \ndos\sin\theta_2$ se escribe como:

\result{Ecuación de Snell 1D (función raíz)}{
\[
f(u) = \nuno\frac{u}{d_1(u)} - \ndos\frac{D-u}{d_2(u)} = 0
\]
con derivada:
\[
f'(u) = \nuno\frac{s_z^2}{d_1^3(u)} + \ndos\frac{t_z^2}{d_2^3(u)} > 0 \quad \text{siempre}
\]
La función $f$ es estrictamente monótona creciente $\Rightarrow$ la raíz es \textbf{única}.
}

La figura~\ref{fig:snell_1d} muestra la función $f(u)$ y la convergencia de Newton-Raphson.

\begin{figure}[H]
\centering
\includegraphics[width=0.82\textwidth]{pv_fig03_snell_1d.pdf}
\caption{Función $f(u) = \nuno u/d_1 - \ndos(D-u)/d_2$ para una geometría típica
($s_z=100\,\text{m}$, $|t_z|=5\,\text{m}$, $D=80\,\text{m}$, $\ndos=2$).
(a) $f(u)$ es monótona creciente con raíz única $u^*\approx 65.6\,\text{m}$.
(b) Las primeras 4 iteraciones de Newton-Raphson (flechas rojas) convergen desde
el estimado inicial $u_0$ con reducción cuadrática del error — típicamente < 5 iteraciones
vs.\ 50-150 de un optimizador quasi-Newton 2D (\texttt{fminunc}).}
\label{fig:snell_1d}
\end{figure}

\section{Algoritmo Newton-Raphson implementado}

El algoritmo implementado en \texttt{calculateRefractionPointFermat.m} es:

\begin{algorithm}[H]
\caption{Punto de refracción de Snell (Newton-Raphson 1D)}
\begin{algorithmic}[1]
\Require Posición fuente $S$, target $T$, índices $\nuno$, $\ndos$
\State $\hat{d} \gets (T_{xy}-S_{xy})/D$,\quad $s_z \gets |S_z|$,\quad $t_z \gets |T_z|$
\State $u \gets D\cdot\dfrac{s_z/\nuno}{s_z/\nuno + t_z/\ndos}$ \Comment{Estimado inicial: interpolación óptica}
\For{$\text{iter} = 1\ldots 12$}
  \State $d_1 \gets \sqrt{u^2+s_z^2}$;\quad $d_2 \gets \sqrt{(D-u)^2+t_z^2}$
  \State $f \gets \nuno\,u/d_1 - \ndos\,(D-u)/d_2$
  \State $f' \gets \nuno\,s_z^2/d_1^3 + \ndos\,t_z^2/d_2^3$
  \State $u \gets \text{clip}(u - f/f',\;0,\;D)$
  \If{$|f/f'| < 10^{-10}\cdot D$} \textbf{break} \EndIf
\EndFor
\State \Return $P^* = [S_{xy} + u\,\hat{d};\;0]$
\end{algorithmic}
\end{algorithm}

\nota{Ganancia de rendimiento}{
Reemplazar \texttt{fminunc} (quasi-Newton 2D, $\sim$100 evaluaciones de la función de costo óptico) por Newton-Raphson 1D (3-5 iteraciones con aritmética elemental) produce una aceleración de \textbf{$\sim$100×} en este cómputo, que es el cuello de botella principal del simulador.
}

%%═══════════════════════════════════════════════════════════════════════════════
\chapter{La Ecuación de Radar Biestático con Refracción}
%%═══════════════════════════════════════════════════════════════════════════════

\section{Propagación biestática en dos medios}

En el sistema biestático, la señal sigue cuatro segmentos de trayectoria con refracción en la interfaz:

\begin{figure}[H]
\centering
\includegraphics[width=0.85\textwidth]{pv_fig04_bistatic_radar.pdf}
\caption{Trayectoria completa de la señal biestática con dos refracciones.
\textbf{Tramo TX→Target}: la señal viaja $R_1=\|TX - Q_1\|$ en aire y $R_2=\|Q_1 - P\|$ en suelo,
refractándose en $Q_1^*$ (punto de Fermat). \textbf{Tramo Target→RX}: análogamente,
$R_3=\|P - Q_2\|$ en suelo y $R_4=\|Q_2 - RX\|$ en aire, refractándose en $Q_2^*$.
Las longitudes totales de trayectoria son $\RT = R_1+R_2$ y $\RR = R_3+R_4$.}
\label{fig:bistatic_radar}
\end{figure}

\section{Derivación de la ecuación de radar con refracción}

La ecuación de radar biestática estándar (sin refracción) para un blanco de sección eficaz $\sigma$ es:
\[
\Prcv = \frac{\Pt\,G_T\,G_R\,\sigma\,\lambda^2}{(4\pi)^3\,R_{TX}^2\,R_{RX}^2}
\]
donde $R_{TX}$ y $R_{RX}$ son las distancias TX-target y target-RX.

Con refracción en la interfaz, deben hacerse tres modificaciones:

\subparagraph{Modificación 1 — Pérdida de propagación con camino óptico total.}
La pérdida de spreading en el tramo TX→Target usa el camino óptico \emph{total}
$\RT = R_1+R_2$ (y no $R_2$ solamente). El frente de onda, aunque viaja a velocidades
diferentes en cada medio, se propaga como si fuera una esfera con radio efectivo $\RT$
al evaluar la potencia en la interfaz de Huygens. Análogamente, en el tramo Target→RX
se usa $\RR = R_3+R_4$.

\subparagraph{Modificación 2 — Transmitancia de potencia.}
En cada cruce de la interfaz se pierde una fracción de la potencia por reflexión:
\begin{itemize}
  \item Tramo TX→Target: transmitancia $T_1(\theta_{i,1})$ en el ángulo de incidencia $\theta_{i,1}$ del TX.
  \item Tramo Target→RX: transmitancia $T_2(\theta_{i,2})$ en el ángulo de incidencia $\theta_{i,2}$ del Target.
\end{itemize}

\result{Ecuación de Radar Biestática con Refracción}{
\[
\boxed{\Prcv = \frac{\Pt\,G_T\,G_R\,\sigma\,T_1\,T_2\,\lambda^2}{(4\pi)^3\,\RT^2\,\RR^2}}
\]
donde:
\begin{align*}
\RT &= R_1 + R_2 = \|TX - Q_1^*\| + \|Q_1^* - P\| \\
\RR &= R_3 + R_4 = \|P - Q_2^*\| + \|Q_2^* - RX\|\\
T_1 &= \frac{\ndos\cos\theta_{t,1}}{\nuno\cos\theta_{i,1}}\,|t_{TM}(\theta_{i,1})|^2 \\
T_2 &= \frac{\nuno\cos\theta_{i,2}}{\ndos\cos\theta_{t,2}}\,|t_{TM}(\theta_{i,2})|^2
\end{align*}
}

\section{El factor $(4\pi)^3$: origen y justificación}

El denominador $(4\pi)^3$ aparece en la ecuación de radar biestático porque la potencia involucra \textbf{tres} etapas de conversión que cada una contribuye un factor $4\pi$:

\begin{enumerate}
  \item \textbf{Pérdida de propagación TX→Target}: el transmisor irradia sobre la esfera de área $4\pi\RT^2$. La densidad de potencia en el target es $\Pt G_T / (4\pi\RT^2)$.
  \item \textbf{Reradiación del target (scatter)}: el target re-emite con densidad de potencia proporcional a $\sigma/(4\pi\RR^2)$ en dirección del RX (sección eficaz $\sigma$ = área efectiva re-irradiadora).
  \item \textbf{Apertura efectiva del RX}: el receptor tiene apertura efectiva $A_{eff} = G_R\lambda^2/(4\pi)$. Multiplicando la densidad de potencia incidente $[\text{W/m}^2]$ por $A_{eff}$ $[\text{m}^2]$ se obtienen los Watts recibidos.
\end{enumerate}

\atencion{Error conceptual corregido}{
Una versión anterior del código usaba \textbf{cuatro} factores $R^2$ independientes:
$(4\pi)^3 R_1^2 R_2^2 R_3^2 R_4^2$, tratando el punto de refracción $Q^*$ como una fuente secundaria que re-emite isótropa\-mente. Esto es físicamente incorrecto: la interfaz es una superficie dieléctrica, no un scatter. La versión corregida usa los caminos totales $\RT^2$ y $\RR^2$.
}

\section{Parámetros del sistema (banda P)}

\begin{table}[H]
\centering
\caption{Parámetros del sistema SAR biestático en banda P para subsuperficie}
\label{tab:system_params}
\begin{tabular}{llrl}
\toprule
Símbolo & Parámetro & Valor & Justificación \\
\midrule
$f_0$ & Frecuencia portadora & $400\,\text{MHz}$ & Banda P — penetración de subsuelo \\
$\lambda$ & Longitud de onda & $0.75\,\text{m}$ & $\lambda = c/f_0$ \\
$B$ & Ancho de banda chirp & $50\,\text{MHz}$ & $[-25,+25]\,\text{MHz}$ alrededor de $f_0$ \\
$\Pt$ & Potencia transmitida & $100\,\text{kW}$ & \\
$G_T = G_R$ & Ganancia antena & $6\,\text{dBi}$ & Patrón real (aper. az. 14°, elev. 50°) \\
$\sigma$ & Sección eficaz target & $1\,\text{m}^2$ & \\
$\varepsilon_r$ & Permitividad relativa suelo & $4$ & Suelo seco típico \\
$\ndos$ & Índice de refracción & $2$ & $\ndos = \sqrt{\varepsilon_r}$ \\
$\thetaB$ & Ángulo de Brewster & $63.43^\circ$ & $\arctan(\ndos/\nuno)$ \\
$z_P$ & Profundidad de targets & $[-10,\,-2]\,\text{m}$ & 64 targets en grid 3D \\
\bottomrule
\end{tabular}
\end{table}

\section{Justificación de la banda P para penetración de subsuelo}

La figura~\ref{fig:pband_attenuation} compara la atenuación electromagnética en suelo
húmedo para distintas frecuencias.

\begin{figure}[H]
\centering
\includegraphics[width=0.82\textwidth]{pv_fig09_alpha_sensitivity.pdf}
\caption{Atenuación en suelo para distintas bandas. A 10\,GHz (banda X), la atenuación
en suelo húmedo excede $25\,\text{dB/m}$, haciendo imposible detectar objetos a más
de 1--2 cm de profundidad. En banda P ($400\,\text{MHz}$), la atenuación es
$\lesssim 0.5\,\text{dB/m}$, permitiendo detectar targets a 5--10 m de profundidad
con presupuesto de enlace favorable.}
\label{fig:pband_attenuation}
\end{figure}

\section{Estimación del link budget}

Con los parámetros de la tabla~\ref{tab:system_params} y geometría típica ($\RT\approx\RR\approx 165\,\text{m}$, $T_1=T_2=1$ en Brewster):

\begin{align*}
\Prcv &= \frac{10^5 \cdot 4 \cdot 4 \cdot 1 \cdot 1 \cdot 1 \cdot (0.75)^2}{(4\pi)^3 \cdot (165)^2 \cdot (165)^2}\\
    &= \frac{10^5 \cdot 16 \cdot 0.5625}{1984.4 \cdot 2.72\times10^4 \cdot 2.72\times10^4}\\
    &\approx \frac{9\times10^5}{1.47\times10^{12}} \approx 6.1\times10^{-7}\,\text{W}
\end{align*}

\[
\Prcv[\text{dBm}] = 10\log_{10}(6.1\times10^{-7}\times1000) \approx -32\,\text{dBm}
\]

\nota{Nota sobre la estimación}{
El valor de $-32\,\text{dBm}$ supone Brewster exacto ($T_1=T_2=1$) y RCS unitario.
En condiciones reales, la potencia recibida será menor por: ángulo no exactamente Brewster,
RCS real del objeto, atenuación en el suelo (no modelada en la ecuación de Friis simple),
y ganancia de antena menor que el máximo direccional.
}

%%═══════════════════════════════════════════════════════════════════════════════
\chapter{La Función de Costo Multi-Objetivo}
%%═══════════════════════════════════════════════════════════════════════════════

\section{El problema de la heterogeneidad dimensional}

El optimizador del plan de vuelo busca simultáneamente:
\begin{itemize}
  \item \textbf{Maximizar} la potencia recibida $\Prcv$ — típicamente en el rango $10^{-8}$–$10^{-12}\,\text{W}$.
  \item \textbf{Minimizar} el área de resolución $\dxy\cdot\dz$ — típicamente en el rango $10^{-2}$–$10^{-4}\,\text{m}^2$.
\end{itemize}

Una combinación lineal directa $w_1\Prcv + w_2\dxy\dz$ es problemática por dos razones:

\begin{enumerate}
  \item \textbf{Dimensiones físicas diferentes}: W vs.\ m$^2$ — la suma no tiene sentido dimensional.
  \item \textbf{Escalas completamente distintas}: $\Prcv\sim 10^{-10}$ frente a $\dxy\dz\sim 10^{-3}$.
    Cualquier peso $w_2$ que haga comparable la resolución con la potencia resultaría
    $w_2 \sim 10^7\,\text{W/m}^2$, sin interpretación física.
\end{enumerate}

\section{Solución: normalización logarítmica}

La física multiplicativa de la ecuación de radar y de las resoluciones SAR sugiere naturalmente el uso de $\log_{10}$:

\begin{itemize}
  \item $\log_{10}(\Prcv)$ para $\Prcv \sim 10^{-10}\,\text{W}$ da un valor $\approx -10$: adimensional, O(10).
  \item $\log_{10}(\dxy\dz)$ para $\dxy\dz \sim 10^{-3}\,\text{m}^2$ da un valor $\approx -3$: adimensional, O(1--5).
\end{itemize}

Propiedades clave que justifican el uso de $\log_{10}$:
\begin{enumerate}
  \item \textbf{Monotonicidad}: $\log_{10}$ es estrictamente creciente $\Rightarrow$ maximizar $\Prcv$ equivale a maximizar $\log_{10}(\Prcv)$.
  \item \textbf{Adimensionalización}: $\log_{10}(x\,[\text{W}])$ pierde las unidades físicas pero conserva toda la información de ordenamiento relativo.
  \item \textbf{Física multiplicativa}: la ecuación de radar es un producto de factores ($G_T$, $T_1$, $\RT^{-2}$, etc.); en escala logarítmica esto se convierte en una suma, lo que la hace compatible con la estructura de la función de costo aditiva.
\end{enumerate}

\section{La función de costo combinada}

\result{Función de Costo Multi-Objetivo Log-Normalizada}{
\[
\boxed{J = -(1-\alphaR)\log_{10}(\Prcv) + \alphaR\log_{10}(\dxy\cdot\dz)}
\]
\begin{itemize}
  \item $\alphaR = 0$: $J = -\log_{10}(\Prcv)$ — \textbf{solo maximizar potencia} (comportamiento original).
  \item $\alphaR = 1$: $J = \log_{10}(\dxy\dz)$ — \textbf{solo minimizar resolución}.
  \item $\alphaR \in (0,1)$: equilibrio multiobjetivo. El valor por defecto es $\alphaR = 0.3$ (70\% potencia, 30\% resolución).
\end{itemize}
El optimizador \texttt{fmincon} \emph{minimiza} $J$, lo cual corresponde a:\\
maximizar $\Prcv$ (primer término) y minimizar $\dxy\dz$ (segundo término).
}

\section{Análisis de los términos en escala logarítmica}

Para la configuración nominal del sistema (tabla~\ref{tab:system_params}):

\begin{table}[H]
\centering
\caption{Órdenes de magnitud de los términos de costo para geometría típica}
\label{tab:cost_magnitudes}
\begin{tabular}{lccc}
\toprule
Cantidad & Valor típico & $\log_{10}$ & Signo en $J$ \\
\midrule
$\Prcv$ & $10^{-10}$–$10^{-8}\,\text{W}$ & $-10$ a $-8$ & $-(1-\alphaR)\times(-10)=+(8\text{ a }10)(1-\alphaR)$ \\
$\dxy$ & $0.01$–$0.1\,\text{m}$ & $-2$ a $-1$ & — \\
$\dz$ & $0.5$–$5\,\text{m}$ & $-0.3$ a $0.7$ & — \\
$\dxy\dz$ & $5\times10^{-3}$–$0.5\,\text{m}^2$ & $-2.3$ a $-0.3$ & $+\alphaR\times(-2)$ \\
\bottomrule
\end{tabular}
\end{table}

Ambos términos son O(1--10), haciendo la ponderación con $\alphaR\in[0,1]$ significativa.

\section{La frontera de Pareto}

La figura~\ref{fig:pareto} ilustra la frontera de Pareto entre potencia y resolución.

\begin{figure}[H]
\centering
\includegraphics[width=0.88\textwidth]{pv_fig07_pareto.pdf}
\caption{Frontera de Pareto en el espacio ($\Prcv$, $\dxy\dz$). Cada punto de la curva corresponde
a un valor de $\alphaR$ y representa una posición óptima distinta del RX. Para $\alphaR=0$ (extremo
inferior-derecho), el RX maximiza potencia. Para $\alphaR=1$ (extremo superior-izquierdo), el RX
minimiza el área de resolución. El punto de equilibrio $\alphaR^* \approx 0.3$ es donde los gradientes
de ambos objetivos en escala logarítmica son iguales.}
\label{fig:pareto}
\end{figure}

\section{Casos degenerados de la función de costo}

La implementación maneja tres casos especiales:

\begin{enumerate}
  \item \textbf{Potencia nula o negativa}: $\log_{10}$ no está definido. Se usa penalización $-300$ (equivalente a $\Prcv = 10^{-300}\,\text{W}$, precio muy alto en el costo).
  \item \textbf{Resolución infinita}: ocurre cuando $\dxy\to\infty$ (separación biestática $\Delta\phi\to 180^\circ$ con target en el eje) o $\dz\to\infty$. Se usa penalización $+10$ en el término de resolución.
  \item \textbf{$\alphaR = 0$}: el término de resolución no se evalúa (por ramas de código), evitando el cómputo innecesario de $\dxy$ y $\dz$.
\end{enumerate}

\section{Derivación del cruce $\alphaR^*$}

El ángulo de cruce $\alphaR^*$ es aquel en que el gradiente del término de potencia iguala al de resolución. En el óptimo de Brewster, la variación conjunta satisface aproximadamente:

\[
(1-\alphaR^*)\,|\Delta\log_{10}\Prcv| = \alphaR^*\,|\Delta\log_{10}(\dxy\dz)|
\]

Para los valores numéricos típicos ($|\Delta\log_{10}\Prcv|\approx 2$, $|\Delta\log_{10}(\dxy\dz)|\approx 0.3$):
\[
\alphaR^* \approx \frac{|\Delta\log_{10}\Prcv|}{|\Delta\log_{10}\Prcv|+|\Delta\log_{10}(\dxy\dz)|}
= \frac{2}{2.3} \approx 0.87
\]

Por encima de $\alphaR^*$, la función de costo es dominada por la resolución; por debajo, por la potencia.

%%═══════════════════════════════════════════════════════════════════════════════
\chapter{La Estrategia de Búsqueda Multi-Inicio}
%%═══════════════════════════════════════════════════════════════════════════════

\section{No convexidad del problema}

La función de costo $J(x_R, y_R)$ como función de la posición horizontal del RX no es convex: puede tener múltiples mínimos locales. Esto se debe a:
\begin{itemize}
  \item El ángulo biestático $\Delta\phi$ depende de la posición relativa TX-RX, y la función $|\cos(\Delta\phi/2)|$ tiene estructura oscilatoria.
  \item Los coeficientes de Fresnel $T_1, T_2$ tienen un máximo pronunciado (pero único) en $\thetaB$.
  \item El patrón de antena introduce dependencia angular adicional.
\end{itemize}

\atencion{Consecuencia práctica}{
Un optimizador de un único inicio puede quedar atrapado en un mínimo local que no corresponde
al ángulo de Brewster. La búsqueda multi-inicio con candidatos heurísticos mitiga este problema.
}

\section{Los cinco candidatos de inicio}

La figura~\ref{fig:multistart} muestra la ubicación geométrica de los 5 candidatos.

\begin{figure}[H]
\centering
\includegraphics[width=0.78\textwidth]{pv_fig06_multistart.pdf}
\caption{Vista cenital (plano $z=0$) mostrando la posición del TX (cuadrado azul), el centroide
del grid de targets (punto rojo), y los 5 candidatos de inicio para la posición del RX.
Los candidatos 1, 2 y 3 cubren las direcciones opuesta y perpendiculares al TX.
Los candidatos 4 y 5 se ubican en la región guiada por el ángulo de Brewster, con jitter
aleatorio para romper simetrías.}
\label{fig:multistart}
\end{figure}

Los candidatos están definidos en \texttt{chooseCandidateRxPos.m} como:

\begin{description}
  \item[Candidato 1 (opuesto):] $\mathbf{x}_0 = \mathbf{c} + d\,\hat{e}_{\text{opp}} \cdot 0.8 + \mathcal{N}(0,1)^2$, donde $\mathbf{c}$ es el centroide del grid de targets, $d$ es la distancia TX-centroide, y $\hat{e}_{\text{opp}}$ apunta desde el centroide en dirección opuesta al TX. Esto aproxima un ángulo biestático $\Delta\phi \approx 180^\circ$.

  \item[Candidato 2 (perpendicular +):] $\mathbf{x}_0 = \mathbf{c} + d\,\hat{e}_{\perp} \cdot 0.8$. Ángulo biestático $\Delta\phi \approx 90^\circ$.

  \item[Candidato 3 (perpendicular $-$):] $\mathbf{x}_0 = \mathbf{c} - d\,\hat{e}_{\perp} \cdot 0.8$. Exploración del semiplano contrario.

  \item[Candidatos 4 y 5 (Brewster):] Se calcula la \textbf{distancia de Brewster}:
  \[
  d_B = \frac{z_{RX} + |z_P|}{\tan\thetaB}
  \]
  y se ubican los candidatos en la dirección azimutal opuesta al TX desde el centroide, a distancia $d_B$, con perturbación aleatoria uniforme de amplitud $0.2\,d_B$.
\end{description}

\section{El optimizador interior-point}

\texttt{fmincon} con algoritmo \texttt{interior-point} minimiza $J(x_R, y_R)$ con restricciones de caja:
\[
x_R \in [-200,\,200]\,\text{m}, \quad y_R \in [-200,\,200]\,\text{m}
\]

Los parámetros de parada son: \texttt{MaxIterations}=200, \texttt{FunctionTolerance}=$10^{-10}$.

En cada llamada a $J$, se evalúan:
\begin{itemize}
  \item Para los 64 targets: 1 llamada Newton-Raphson (Target→RX) + cálculos de Fresnel/radar.
  \item 1 llamada a \texttt{calculateBistaticResolution} (aritmética pura, sin optimización interna).
\end{itemize}

\section{Pre-cómputo del tramo TX→Target}

Una optimización clave es que el tramo TX→Target (refracción $Q_1^*$, ángulo $\theta_{i,1}$, transmitancia $T_1$, longitud $\RT$) es \textbf{constante} para todas las posiciones candidatas del RX, ya que TX y los targets no cambian durante la optimización. Por tanto, estos datos se pre-computan una sola vez por posición TX (función \texttt{precomputeTxData.m}) antes del bucle multi-inicio:

\begin{figure}[H]
\centering
\includegraphics[width=0.92\textwidth]{pv_fig05_cost_function.pdf}
\caption{Función de costo $J(x_R, y_R)$ en el plano horizontal para $\alphaR=0.3$.
(a) Mapa de calor de $J$ para una posición TX fija, mostrando la estructura de mínimos múltiples.
(b) Perfil de $J$ a lo largo de la dirección azimutal óptima: el mínimo global corresponde
a la posición de Brewster opuesta al TX. Los 5 candidatos de inicio se marcan con cruces.}
\label{fig:cost_function}
\end{figure}

%%═══════════════════════════════════════════════════════════════════════════════
\chapter{Estimaciones Numéricas y Análisis de Sensibilidad}
%%═══════════════════════════════════════════════════════════════════════════════

\section{Geometría nominal y parámetros derivados}

Con los parámetros de la tabla~\ref{tab:system_params} y la hélice cónica nominal:

\begin{table}[H]
\centering
\caption{Parámetros geométricos de la hélice cónica y distancias derivadas}
\label{tab:geometry}
\begin{tabular}{llrl}
\toprule
Símbolo & Parámetro & Valor & Cálculo \\
\midrule
$\rho_{min}$ & Radio menor hélice & $147.5\,\text{m}$ & — \\
$\rho_{max}$ & Radio mayor hélice & $172.5\,\text{m}$ & — \\
$z_{min}$ & Altitud mínima (base) & $80\,\text{m}$ & — \\
$z_{max}$ & Altitud máxima (cima) & $120\,\text{m}$ & — \\
$\Delta\rho$ & Variación de radio & $25\,\text{m}$ & $\rho_{max}-\rho_{min}$ \\
$\Delta z$ & Variación de altura & $40\,\text{m}$ & $z_{max}-z_{min}$ \\
$B_{helix}$ & Longitud espiral & $47.17\,\text{m}$ & $\sqrt{\Delta\rho^2+\Delta z^2}$ \\
$\beta$ & Ángulo inclinación & $58.0^\circ$ & $\arctan(\Delta z/\Delta\rho)$ \\
$\rho_0$ & Radio medio & $160\,\text{m}$ & $(\rho_{min}+\rho_{max})/2$ \\
$z_0$ & Altitud media & $100\,\text{m}$ & $(z_{min}+z_{max})/2$ \\
$\psi_0$ & Look angle medio & $58.0^\circ$ & $\arctan(\rho_0/z_0)$ \\
$d_B$ & Distancia Brewster & $\approx 105\,\text{m}$ & $(z_0+|z_P|)/\tan\thetaB$, $z_P=-5\,\text{m}$ \\
\bottomrule
\end{tabular}
\end{table}

Se verifica que $\beta \approx \psi_0 \approx 58^\circ$, lo que significa que la hélice tiene inclinación óptima ($B_\perp \approx B_{helix}$) para la resolución vertical. Esto \emph{no es coincidencia} — la hélice fue diseñada con esta condición.

\section{Estimación de la resolución 3D en el punto óptimo}

Aplicando las fórmulas de \texttt{explicacion\_resolucion.tex} con $\Delta\phi=90^\circ$ (NorthOffset del RX):

\subparagraph{Ángulo de transmisión en el suelo:}
\[
\cos\theta_{t,0} = \sqrt{1-\frac{\sin^2(58^\circ)}{4}} = \sqrt{1-0.181} = 0.905
\]

\subparagraph{Distancia de slant range:}
\[
R_0 = \sqrt{\rho_0^2 + z_0^2} = \sqrt{160^2 + 100^2} = 188.7\,\text{m}
\]

\subparagraph{Apertura tomográfica efectiva:}
\[
B_\perp = B_{helix}\,|\cos(\beta-\psi_0)| = 47.17\,|\cos(0^\circ)| = 47.17\,\text{m}
\]

\subparagraph{Resolución vertical:}
\begin{align*}
W_z &= \ndos\,B\,\cos\theta_{t,0} + \frac{c\,B_\perp\,\sin\psi_0\cos\psi_0}{\lambda\,R_0\,\ndos\cos\theta_{t,0}}\\
    &= 2 \times 50\times10^6 \times 0.905 + \frac{3\times10^8 \times 47.17 \times \sin(58^\circ)\cos(58^\circ)}{0.75 \times 188.7 \times 2 \times 0.905}\\
    &= 90.5\,\text{MHz} + \frac{3\times10^8 \times 47.17 \times 0.449}{256.2}\\
    &= 90.5\,\text{MHz} + 24.9\,\text{MHz} = 115.4\,\text{MHz}
\end{align*}
\[
\dz = \frac{c}{2\,W_z} = \frac{3\times10^8}{2 \times 115.4\times10^6} = 1.30\,\text{m}
\]

\subparagraph{Resolución horizontal ($\Delta\phi=90^\circ$):}
\[
\dxy = \frac{0.60\,\lambda}{\pi\,\sin\psi_0\,|\cos(\Delta\phi/2)|}
= \frac{0.60 \times 0.75}{\pi \times \sin(58^\circ) \times |\cos(45^\circ)|}
= \frac{0.45}{\pi \times 0.848 \times 0.707} = 0.24\,\text{m}
\]

\result{Resoluciones nominales estimadas (P-band, $\Delta\phi=90^\circ$)}{
\begin{align*}
\dxy &\approx 0.24\,\text{m} \quad \text{(horizontal)} \\
\dz &\approx 1.30\,\text{m} \quad \text{(vertical)} \\
\dxy\cdot\dz &\approx 0.31\,\text{m}^2
\end{align*}
La contribución del ancho de banda al $W_z$ es 90.5 MHz (78\%) y la contribución
tomográfica es 24.9 MHz (22\%); el término de BW domina en banda P.
}

\section{Análisis de sensibilidad de $J$ respecto a $\alphaR$}

La figura~\ref{fig:alpha_sensitivity} muestra cómo la posición óptima del RX y el valor de $J^*$ cambian con $\alphaR$.

\begin{figure}[H]
\centering
\includegraphics[width=0.90\textwidth]{pv_fig09_alpha_sensitivity.pdf}
\caption{Sensibilidad de la optimización al parámetro $\alphaR$.
(a) Valor del costo $J^*$ en el óptimo como función de $\alphaR$: el mínimo se desplaza
de la región de Brewster ($\alphaR\to 0$) hacia la región de mejor resolución ($\alphaR\to 1$).
(b) Componentes individuales: contribución de potencia (azul) y de resolución (rojo)
en el punto óptimo para cada $\alphaR$.
El cruce $\alphaR^*\approx 0.87$ indica el punto donde ambos objetivos tienen igual peso
en la función de costo total.}
\label{fig:alpha_sensitivity}
\end{figure}

\section{Vista general del sistema optimizado}

La figura~\ref{fig:sistema_3d} muestra la configuración completa del sistema con la
trayectoria del RX optimizada para $\alphaR=0.3$.

\begin{figure}[H]
\centering
\includegraphics[width=0.90\textwidth]{pv_fig10_sistema_3d.pdf}
\caption{Vista 3D del sistema optimizado. (a) La hélice cónica del TX (línea azul sólida)
desciende mientras amplía su radio. (b) Las posiciones óptimas del RX (puntos rojos)
forman una trayectoria aproximada que mantiene el ángulo biestático cercano al de Brewster.
(c) Las flechas naranjas ilustran los caminos refractados TX→Target y Target→RX para la
posición óptima de un TX seleccionado.}
\label{fig:sistema_3d}
\end{figure}

%%═══════════════════════════════════════════════════════════════════════════════
\chapter{Conclusiones y Guía de Uso}
%%═══════════════════════════════════════════════════════════════════════════════

\section{Resumen del método}

\result{Algoritmo de Optimización del Plan de Vuelo (resumen)}{
\textbf{Para cada posición TX en la hélice}:
\begin{enumerate}
  \item Pre-computar ganancias TX $G_T^{(i)}$ y datos de Fermat TX→Target ($R_T^{(i)}, T_1^{(i)}$) para los 64 targets.
  \item Construir la geometría de búsqueda (distancia de Brewster, dirección opuesta).
  \item Para cada uno de los 5 candidatos de inicio:
    \begin{enumerate}
      \item Inicializar $\mathbf{x}_0^{(k)}$ según las heurísticas geométricas.
      \item Minimizar $J(\mathbf{x}_R)$ con \texttt{fmincon} (bounds $\pm 200\,\text{m}$).
      \item En cada evaluación de $J$: resolver Newton-Raphson 1D para Target→RX (64×), calcular $\dxy\cdot\dz$.
    \end{enumerate}
  \item Seleccionar el candidato con menor $J^*$.
  \item Reportar $\Prcv^* = 10^{-J_\text{power}^*}$ y las resoluciones $\dxy^*, \dz^*$ en el óptimo.
\end{enumerate}
}

\section{Guía de parámetros de diseño}

\begin{enumerate}
  \item \textbf{Selección de $\alphaR$}: comenzar con $\alphaR=0$ (solo potencia) para verificar que el ángulo de Brewster se encuentra. Luego incrementar hacia $\alphaR=0.3$. No superar $\alphaR^*\approx 0.87$ sin verificar que la posición de Brewster siga siendo accesible.

  \item \textbf{Número de candidatos de inicio}: 5 candidatos es el mínimo razonable para cubrir los cuadrantes de búsqueda. Incrementar a 8--10 si el espacio es muy no convexo (e.g., grid de targets muy ancho).

  \item \textbf{Factor de decimación}: el código decimóa la trayectoria TX en un factor 200 (por defecto), lo que produce $\sim\!25$ posiciones TX evaluadas por vuelta. Reducir a 100 para mayor resolución temporal de la trayectoria del RX; aumentar a 400+ si el tiempo de cómputo es limitante.

  \item \textbf{Frecuencia de operación}: mantener en banda P ($\leq 500\,\text{MHz}$) para targets a más de 1 m de profundidad. La penalización en resolución horizontal ($\dxy\propto\lambda$) es severa pero inevitable para la penetración requerida.

  \item \textbf{Condición $\beta\approx\psi_0$}: verificar que la inclinación de la hélice satisfaga $\beta = \arctan(\Delta z/\Delta\rho) \approx \arctan(\rho_0/z_0)$. Esto maximiza $B_\perp$ y por ende $W_z$, mejorando la resolución vertical sin cambiar la trayectoria TX.
\end{enumerate}

\section{Limitaciones}

\begin{itemize}
  \item La ecuación de radar no incluye \textbf{atenuación dieléctrica} en el suelo ($e^{-\alpha|z_P|}$). Para subsuelos húmedos, este término reduce significativamente la potencia recibida.
  \item Los coeficientes de Fresnel asumen \textbf{medios homogéneos sin pérdidas} ($\varepsilon_r$ real). En suelos reales, $\varepsilon_r$ es complejo y el ángulo de Brewster se modifica.
  \item El grid de 64 targets es \textbf{representativo} del volumen de interés, no una discre\-tización exhaustiva. Los targets reales son continuos.
  \item La \textbf{interpolación PCHIP} de la trayectoria del RX entre posiciones optimizadas puede introducir oscilaciones si el decimationFactor es demasiado grande.
\end{itemize}

%%═══════════════════════════════════════════════════════════════════════════════
\appendix
\chapter{Derivación Formal de la Reducción 3D$\to$1D de Fermat}
%%═══════════════════════════════════════════════════════════════════════════════

\section{Enunciado del teorema}

\begin{mdframed}
\textbf{Teorema (Fermat 1D para interfaz plana).}
Sea $S = (s_x, s_y, s_z)$ con $s_z > 0$ una fuente en el Medio 1 ($n_1$) y $T = (t_x, t_y, t_z)$ con $t_z < 0$ un target en el Medio 2 ($n_2$). El punto de refracción $P^* = \arg\min_{P\in z=0} [n_1\|S-P\|+n_2\|P-T\|]$ pertenece al segmento de línea horizontal que une $S_{xy}=(s_x,s_y)$ con $T_{xy}=(t_x,t_y)$ proyectados en $z=0$.
\end{mdframed}

\section{Demostración}

Sea $P = (p_x, p_y, 0)$ y definamos $\mathcal{T}(p_x,p_y) = n_1 d_1 + n_2 d_2$ donde:
\[
d_1 = \sqrt{(p_x-s_x)^2+(p_y-s_y)^2+s_z^2}, \quad d_2 = \sqrt{(p_x-t_x)^2+(p_y-t_y)^2+t_z^2}
\]

Las condiciones de primer orden $\partial\mathcal{T}/\partial p_x = \partial\mathcal{T}/\partial p_y = 0$ dan:
\[
n_1\frac{p_x-s_x}{d_1} + n_2\frac{p_x-t_x}{d_2} = 0
\]
\[
n_1\frac{p_y-s_y}{d_1} + n_2\frac{p_y-t_y}{d_2} = 0
\]

Estas dos ecuaciones pueden escribirse vectorialmente como:
\[
n_1\frac{\mathbf{p}-\mathbf{S}_{xy}}{d_1} = -n_2\frac{\mathbf{p}-\mathbf{T}_{xy}}{d_2}
\]

donde $\mathbf{S}_{xy}=(s_x,s_y)$, $\mathbf{T}_{xy}=(t_x,t_y)$, $\mathbf{p}=(p_x,p_y)$.

Denotando $\mathbf{a} = \mathbf{p}-\mathbf{S}_{xy}$ y $\mathbf{b} = \mathbf{T}_{xy}-\mathbf{p}$, la condición se convierte en:
\[
\frac{n_1}{d_1}\mathbf{a} = \frac{n_2}{d_2}\mathbf{b}
\]

Esto significa que $\mathbf{a}$ y $\mathbf{b}$ son \textbf{paralelos} (proporcionales por un escalar positivo). Geométricamente, $\mathbf{a} = \mathbf{p}-\mathbf{S}_{xy}$ y $\mathbf{b} = \mathbf{T}_{xy}-\mathbf{p}$ paralelos implican que $\mathbf{p}$, $\mathbf{S}_{xy}$ y $\mathbf{T}_{xy}$ son colineales, es decir, $\mathbf{p}$ pertenece a la línea que pasa por $\mathbf{S}_{xy}$ y $\mathbf{T}_{xy}$. $\square$

\section{Unicidad de la raíz}

Parametrizando $\mathbf{p} = \mathbf{S}_{xy} + u\hat{d}$ con $D = \|\mathbf{T}_{xy}-\mathbf{S}_{xy}\|$ y $\hat{d} = (\mathbf{T}_{xy}-\mathbf{S}_{xy})/D$, la función:
\[
f(u) = n_1\frac{u}{\sqrt{u^2+s_z^2}} - n_2\frac{D-u}{\sqrt{(D-u)^2+t_z^2}}
\]
tiene derivada:
\[
f'(u) = \frac{n_1\,s_z^2}{(u^2+s_z^2)^{3/2}} + \frac{n_2\,t_z^2}{((D-u)^2+t_z^2)^{3/2}} > 0
\]
para todos $u\in\mathbb{R}$, $s_z>0$, $t_z\neq 0$. Por tanto $f$ es estrictamente monótona creciente, tiene a lo sumo una raíz. Combinando con los valores de frontera $f(0) = -n_2\,D/\sqrt{D^2+t_z^2} < 0$ y $f(D) = n_1\,D/\sqrt{D^2+s_z^2} > 0$, el teorema del valor intermedio garantiza exactamente una raíz $u^*\in(0,D)$. $\square$

%%═══════════════════════════════════════════════════════════════════════════════
\backmatter
\chapter*{Referencias}
\addcontentsline{toc}{chapter}{Referencias}
%%═══════════════════════════════════════════════════════════════════════════════

\begin{enumerate}
  \item M.\ Born y E.\ Wolf, \emph{Principles of Optics}, 7ª ed., Cambridge University Press, 1999. [Ecuaciones de Fresnel, ángulo de Brewster]
  \item D.\ C.\ Munson, J.\ D.\ O'Brien, W.\ K.\ Jenkins, ``A tomographic formulation of spotlight-mode SAR,'' \emph{Proc.\ IEEE}, vol.\ 71, no.\ 8, 1983.
  \item J.\ A.\ Góes, \emph{Techniques for High-Resolution 3D Images with Synthetic Aperture Radar}, Tesis Doctoral, UNICAMP, 2022. [Resoluciones $\dxy$, $\dz$ para SAR biestático helicoidal]
  \item F.\ T.\ Ulaby, R.\ K.\ Moore, A.\ K.\ Fung, \emph{Microwave Remote Sensing}, Artech House, 1981. [Atenuación en suelo, permitividades]
  \item G.\ F.\ Edelmann et al., ``Imaging subsurface objects with P-band SAR,'' \emph{IEEE TGRS}, 2019. [Motivación banda P]
  \item R.\ Kress, \emph{Numerical Analysis}, Springer, 1998. [Convergencia de Newton-Raphson]
  \item A.\ Ishimaru, \emph{Electromagnetic Wave Propagation, Radiation, and Scattering}, Prentice Hall, 1991. [Ecuación de radar biestática]
  \item O.\ Arikan, D.\ C.\ Munson Jr., ``A tomographic formulation of bistatic SAR,'' \emph{ICASSP}, 1988.
\end{enumerate}

\end{document}
